% Letter 'murica
% 2.23in margins all around gives you the golden ratio between content width
% margins apparently
\documentclass[letterpaper, 12pt]{article}
\usepackage[letterpaper,
	top=1in,
	bottom=1.25in,
	left=1.5in,
	right=1.5in,
]{geometry}

\usepackage[T1]{fontenc}
\usepackage{ETbb}
\let\oldstylenums\textosf

%% Exclusively use indents to demarcate paragraphs
\setlength{\parindent}{0.5in}
%% Don't add any more space than baselineskip after paragraphs
\setlength{\parskip}{0in}

% 135% of the font size is the ideal line spacing
% 1.35 * 12 = 16.20
% 16.20/14.5 (default baselineskip) = 1.116
% \linespread scales \baselineskip ergo the math shenanigans
\linespread{1.117}

\RequirePackage[usenames,dvipsnames]{color}
\definecolor{uvaOrange}{RGB}{229, 114, 0}

\usepackage{fancyhdr}
\pagestyle{fancy}
\renewcommand{\headrulewidth}{0pt}
\renewcommand{\footrulewidth}{0pt}
\fancyhead{}
\fancyfoot{}
\fancyfoot[C]{{\color{uvaOrange}\thepage}}

\usepackage[hidelinks, colorlinks]{hyperref}
% END OF GENERIC DBK_ARTICLE

\usepackage[style=alphabetic]{biblatex}
\addbibresource{ref.bib}

\usepackage{minted}
\usepackage[title]{appendix}
\usepackage{booktabs}
\usepackage{tabularx}
\usepackage{float}

\title{Hypothesis}
\author{Deebakkarthi Chinnasame Rani}

\usepackage{lipsum}
\begin{document}

%% Title
\maketitle
\thispagestyle{fancy}

%%
\begin{abstract}
\end{abstract}

\section{Tiered Verification}

\section{What is Verification}
Merriam Webster defines \textbf{verify} as \textit{to confirm or substantiate
in law by oath}.
In formal hardware verification, we are confirming that our implementation
does everything that we had designed (and if we're lucky, \textbf{only} what
we designed). This process fundamentally requires a ground truth.

In the absence of such a ground truth, the properties generated by the model
can only be considered as \textit{candidates}. They are a guess on part of the
model on what the designer intended in their RTL.
\cite{Foster2005}

\section{Verilog Knowledge Benchmark}
Drawing inspiration from \cite{phanBenchmarkExpertlevelAcademic2026}, I created
a benchmark to test the Verilog Knowledge of a LLM. To automate the manual
version I had last year, I split the book
\citetitle{chonnadVerilogFrequentlyAsked2004}
\cite{chonnadVerilogFrequentlyAsked2004}
on chapter boundaries. Each chapter was then processed using the
prompts described in Appendix \ref{sec:a_prompts} and converted into a
markdown document of questions.

\subsection{Results}
The models have much better Verilog knowledge that they used to. As we see in
Table \ref{tab:seq_results}, \texttt{opus} nearly aces this whereas last year
I was able to fool the frontier model with minimal chichanery.

\begin{table}[]
	\centering
	\begin{tabular}{@{}lll@{}}
		\toprule
		Model  & Correct & Accuracy \\ \midrule
		\texttt{opus}   & 216/221 & 97.7\%   \\ \midrule
		\texttt{sonnet} & 211/221 & 95.5\%   \\ \midrule
		\texttt{haiku}  & 196/221 & 88.7\%   \\ \bottomrule
	\end{tabular}
	\caption{Sequential Results}
	\label{tab:seq_results}
\end{table}

\subsection{Sequential vs Batch Mode}
The results shown in Table \ref{tab:seq_results} are of the \textit{sequential}
benchmark. Each question was answered in isolation. This meant 221 individual
API call and each run taking about 10 minutes to complete.

I attempted to hasten this process by \textit{batching} them in to a
single prompt. I figured this shouldn't have that big of an impact as the entire
JSON file was just 25000 words, a miniscule amount compared to the 1M token
context window. But I discovered significant intelligence drop off, even in
\texttt{opus}. Table \ref{tab:batch_results} illustrates this. We see that
even sequential \texttt{haiku} is able to beat batch-mode \texttt{opus}!

\begin{table}[H]
	\centering
	\begin{tabular}{@{}lll@{}}
		\toprule
		Model  & Correct & Accuracy \\ \midrule
		opus   & 156/221 & 70.6\%   \\ \midrule
		sonnet & 128/221 & 57.9\%   \\ \midrule
		haiku  & 98/221  & 44.3\%   \\ \bottomrule
	\end{tabular}
	\caption{Batch Results}
	\label{tab:batch_results}
\end{table}

Furthermore, \textbf{none} of the model answered all the questions.
\texttt{opus} and \texttt{haiku} truncated their responses to around 200
questions whereas \texttt{sonnet} returned 269 lines.

\subsection{Takeaways}
This further leads to me to lean towards a \textit{globalized planning-localized
generation} approach where we first prompt the LLM , with the entire context, to
perform division of labour and then spawn separate agents for each task.

Concretely, this would look like using individual modules instead of entire
RTL files or, at a more granular level, using static analysis to single out
lexical scope instead of entire modules. To capture inter-scope/module
dependencies a separate prompt would be made.

\section{Don't use SystemVerilog}
Based on the numerous benchmark runs I've done, the vastness of
SystemVerilog syntax bogs down the model. I'm certain that the frontier
models can easily fix the compiler errors when allowed to iterate,
but I'm hesitant
to do so as this could all be avoided by just not using SystemVerilog.

The \textbf{only} task that necessitates an LLM's expertise is the extraction
of temporal properties. We can concisely express them using LTL's terse syntax.
All other tasks are merely plumbing to make this flow automated.
Though it may be
arduous, all these task can be done entirely through static analysis.
This saves us time and tokens.

I have also noticed that the properties extracted are different when the LLM
operates purely on LTL compared to when it has to generate compilable
SystemVerilog. Generating LTL also permits us to exploit the plethora of
LTL tools to perform redundancy checks, triviality checks, etc.

\printbibliography

\begin{appendices}

	\section{Prompts}
	\label{sec:a_prompts}

\begin{minted}[breaklines]{markdown}
# Agent Instructions: Convert Verilog FAQ Chapter to HLE-Style Questions

## Your Task

You are converting one chapter of the book **"Verilog: Frequently Asked Questions — Language, Applications and Extensions"** (Chonnad & Balachander, Springer 2004) into benchmark questions modelled on the **Humanity's Last Exam (HLE)** style.

**THIS PHASE ONLY: Convert existing questions. Do NOT invent new questions.**
Every question you produce must be directly grounded in a Q&A that already exists in the chapter. You are reformatting and sharpening, not authoring.

---

## Source Material

You will be given a path to a chapter PDF. Read it in full before writing any questions.
If you want a plain-text summary of a section or a page, **spawn a sub-agent** and ask it to read and summarise the relevant pages for you — do not guess at content you haven't read.

The book is structured as numbered FAQ entries, e.g.:
> **1.1.3** What are the differences between blocking and nonblocking assignments?

Each entry has:
- A question headline
- A prose explanation (often with Verilog code examples)
- Sometimes a comparison table or waveform description

Your raw material is those entries.

---

## HLE Question Style (summary)

Full style guide is at `/Users/deebakkarthi/.local/src/faq_extraction/QUESTION_FORMAT.md`. Key rules:

1. **Dense technical stem** — Use precise Verilog/HDL terminology. Do not simplify or paraphrase into plain English. The question should be hard to parse without domain knowledge.
2. **Single unambiguous answer** — Every question has exactly one correct answer, verifiable by a Verilog expert.
3. **Short exact answer** — Either one letter (MCQ) or a short exact string (short-answer).
4. **MCQ: 5 adversarially close options (A–E)** — All options must be in the same concept-family. Wrong options should be things a model reasoning incorrectly would plausibly choose.
5. **Non-trivial** — The answer must not be obvious to anyone who merely glanced at the question. It should require knowing the specific Verilog rule or behaviour.
6. **No open-ended answers** — No "explain why", "describe", "give an example". Closed-ended only.
7. **No artificial difficulty** — Do not ask questions whose difficulty comes from requiring a simulator or mental execution of long code. The difficulty must come from knowledge of the language rule.

### What makes a good Verilog MCQ distractor?

For a question about a Verilog rule, wrong options should be:
- The behaviour in a *similar but different* construct (e.g., blocking vs. nonblocking, `initial` vs. `always`)
- The behaviour under a *different simulator phase* or event ordering
- A common misconception explicitly mentioned in the book
- The correct behaviour in a *related but distinct* scenario

Avoid distractors that are obviously wrong to anyone with basic Verilog knowledge.

---

## Output Format

Write your output to the file specified by the caller (e.g., `CHAPTER_01.md`).

Each question entry must follow this exact structure:

```
### Q{chapter}{index} — {source_section}

**Stem:**
{Full question text, including any setup. Use code blocks for Verilog snippets.}

**Answer Choices:**
A. {option}
B. {option}
C. {option}
D. {option}
E. {option}

**Answer:** {letter}

**Source:** Section {X.Y.Z} — "{original FAQ question headline}"

**Rationale:** {One or two sentences explaining why the correct answer is right and why the key distractors are wrong.}

---
```

For short-answer questions (use sparingly, only when MCQ would be awkward):

```
### Q{chapter}{index} — {source_section}

**Stem:**
{Full question text. End with an explicit format instruction, e.g. "Answer with a keyword."}

**Answer:** {exact string}

**Source:** Section {X.Y.Z} — "{original FAQ question headline}"

**Rationale:** {One or two sentences.}

---
```

Use `chapter` = two-digit chapter number (01, 02, …) and `index` = three-digit sequential index (001, 002, …).

---

## Self-Containment Rule (CRITICAL)

**Every question must be fully self-contained. A reader with no access to this book must be able to answer it.**

This means:
- **Never reference the book, a table, a figure, or a section by name.** Do not write "According to Table 1-2", "as shown in the book", "the text states", "in the example above", "per Figure 3-1", etc. The question must stand alone.
- **If a table is the source of a fact, reproduce the relevant row(s) inline in the stem** — or restate the rule as a self-contained technical statement.
- **If a code example is the source, include the code inline** in a Verilog code block within the stem.
- **If a rule comes from the book, state it as a Verilog language rule**, not as a claim the book makes. Write "In Verilog, X is true" not "The book states X is true".
- Any context the reader needs to answer the question must appear inside the `**Stem:**` field itself.

Violations of this rule make the question unanswerable without the book and disqualify it.

---

## Quality Bar

Before writing a question, ask yourself:
- Could a Verilog engineer who has never seen this book answer this from the stem alone? (If not, embed the missing context or discard)
- Is there any ambiguity about which answer is correct? (Bad — tighten the stem or discard)
- Are all 5 MCQ options things a partially-knowledgeable engineer might genuinely consider? (If one option is obviously absurd, replace it)
- Does the answer appear verbatim in the question stem? (Bad — rephrase)
- Does the stem reference "the book", "the table", "the figure", or "the example"? (Bad — rewrite to embed the context)

## Coverage

- Convert **every** FAQ entry in the chapter that can be expressed as a closed-ended question.
- Skip entries that are purely illustrative (e.g., "here is an example of X") with no testable fact.
- Skip entries whose answer is inherently open-ended or simulator-dependent.
- Aim for one question per FAQ entry. If an entry contains multiple distinct testable facts, produce one question per fact.

---

## Checklist Before Saving

- [ ] Read the full chapter PDF before writing any questions
- [ ] Every question traces to a specific section number
- [ ] No question is invented — all grounded in existing book content
- [ ] All MCQ options are in the same concept-family
- [ ] Correct answer is unambiguous
- [ ] No open-ended questions
- [ ] File saved to the correct path
\end{minted}

\end{appendices}

\end{document}
